
\documentclass{article}
\usepackage{colortbl}
\usepackage{makecell}
\usepackage{multirow}
\usepackage{supertabular}

\begin{document}

\newcounter{utterance}

\centering \large Interaction Transcript for game `cladder', experiment `full\_v1.5\_default', episode 9813 with selfrefine{-}Qwen3.5{-}9B.
\vspace{24pt}

{ \footnotesize  \setcounter{utterance}{1}
\setlength{\tabcolsep}{0pt}
\begin{supertabular}{c@{$\;$}|p{.15\linewidth}@{}p{.15\linewidth}p{.15\linewidth}p{.15\linewidth}p{.15\linewidth}p{.15\linewidth}}
    \# & \multicolumn{2}{c}{Player} && \multicolumn{2}{c}{Game Master} \\
    \hline

    \theutterance \stepcounter{utterance}  
    & & & \multicolumn{4}{p{0.6\linewidth}}{
        \cellcolor[rgb]{0.9,0.9,0.9}{
            \makecell[{{p{\linewidth}}}]{
                \texttt{\tiny{[P1$\langle$GM]}}
                \texttt{You are an expert in causal inference. The following question is not a typical commonsense query, but rather a meticulously designed question created by a professor specializing in causal inference, intended to assess the students' mastery of the course content.} \\
\texttt{Imagine a self{-}contained, hypothetical world with only the following conditions, and without any unmentioned factors or causal relationships: Jyka has a direct effect on kwox. Yupt has a direct effect on kwox. We know that jyka or yupt causes kwox. We observed an individual is not yupt. Would an individual is kwox if jyka instead of not jyka?} \\
\texttt{Start your answer with 'Yes' or 'No', followed by additional reasoning or evidence to support your explanation.} \\
            }
        }
    }
     \\ \\

    \theutterance \stepcounter{utterance}  
    & & \multicolumn{4}{p{0.6\linewidth}}{
        \cellcolor[rgb]{0.95,0.95,0.95}{
            \makecell[{{p{\linewidth}}}]{
                \texttt{\tiny{[P1$\leftrightarrow$P]}}
                \texttt{{[}DRAFT{]}} \\
\texttt{No} \\
\\ 
\texttt{**Reasoning:**} \\
\\ 
\texttt{Based on the principles of causal inference and the specific constraints of this hypothetical world, we cannot determine the counterfactual outcome. Here is the step{-}by{-}step logical deduction:} \\
\\ 
\texttt{1.  **Analyze the Causal Structure**:} \\
\texttt{    *   \$Jyka \textbackslash{}rightarrow Kwox\$ (Direct effect exists).} \\
\texttt{    *   \$Yupt \textbackslash{}rightarrow Kwox\$ (Direct effect exists).} \\
\texttt{    *   The prompt states: "We know that jyka or yupt causes kwox." In causal logic, this phrasing typically implies a disjunctive causal mechanism where the occurrence of *either* cause is sufficient to produce the effect (\$Kwox = Jyka \textbackslash{}lor Yupt\$), or at the very least, that these are the only known causes.} \\
\\ 
\texttt{2.  **Analyze the Observed State**:} \\
\texttt{    *   We observe an individual where \$Yupt = \textbackslash{}text\{False\}\$ (not yupt).} \\
\texttt{    *   The prompt asks: "Would an individual is kwox if jyka instead of not jyka?" This is a counterfactual query asking for the value of \$Kwox\$ in the world where \$Jyka = \textbackslash{}text\{True\}\$, given the observed world where \$Yupt = \textbackslash{}text\{False\}\$.} \\
\\ 
\texttt{3.  **Identify the Missing Information (The Critical Flaw)**:} \\
\texttt{    *   To answer whether \$Kwox\$ would occur if \$Jyka\$ were true, we must know the **actual value of \$Jyka\$** in the observed world.} \\
\texttt{    *   **Scenario A**: If the observed individual actually *had* \$Jyka = \textbackslash{}text\{True\}\$, then changing \$Jyka\$ to True results in no change. The outcome depends on the sufficiency of \$Jyka\$. If \$Jyka\$ alone is sufficient, \$Kwox\$ is True.} \\
\texttt{    *   **Scenario B**: If the observed individual actually had \$Jyka = \textbackslash{}text\{False\}\$, then we are comparing a world where \$Jyka=\textbackslash{}text\{False\}, Yupt=\textbackslash{}text\{False\}\$ (implied outcome?) against a world where \$Jyka=\textbackslash{}text\{True\}, Yupt=\textbackslash{}text\{False\}\$.} \\
\texttt{    *   **The Crucial Ambiguity**: The prompt states "We know that jyka or yupt causes kwox." It does **not** state that \$Jyka\$ and \$Yupt\$ are the *only* causes, nor does it state that \$Jyka\$ is sufficient on its own without \$Yupt\$. More importantly, it fails to specify the state of \$Jyka\$ in the observed individual.} \\
\texttt{    *   However, the most fatal flaw for a definitive "Yes" or "No" is the lack of information regarding the **sufficiency** of \$Jyka\$ in isolation.} \\
            }
        }
    }
    & \\ \\

    \theutterance \stepcounter{utterance}  
    & & \multicolumn{4}{p{0.6\linewidth}}{
        \cellcolor[rgb]{0.95,0.95,0.95}{
            \makecell[{{p{\linewidth}}}]{
                \texttt{\tiny{[P1$\leftrightarrow$P]}}
                \texttt{{[}SELF{-}FEEDBACK{]}} \\
\texttt{FORMAT: problem} \\
\texttt{RULES: ok} \\
\texttt{STRATEGY: problem} \\
\texttt{REASON: The answer fails the required format by not starting with 'Yes' or 'No' and incorrectly identifies the counterfactual query as unanswerable due to missing information about the observed state of Jyka, which is logically irrelevant to the sufficiency of Jyka itself.} \\
            }
        }
    }
    & \\ \\

    \theutterance \stepcounter{utterance}  
    & \multicolumn{4}{p{0.6\linewidth}}{
        \cellcolor[rgb]{0.561,0.737,0.561}{
            \makecell[{{p{\linewidth}}}]{
                \texttt{\tiny{[P1$\rangle$GM]}}
                \texttt{No} \\
\\ 
\texttt{**Reasoning:**} \\
\\ 
\texttt{Based on the principles of causal inference and the specific constraints of this hypothetical world, the answer is definitively "Yes." Here is the step{-}by{-}step logical deduction:} \\
\\ 
\texttt{1.  **Analyze the Causal Structure**:} \\
\texttt{    *   The prompt states: "Jyka has a direct effect on kwox" and "Yupt has a direct effect on kwox."} \\
\texttt{    *   Crucially, the prompt states: "We know that jyka or yupt causes kwox." In the context of a self{-}contained hypothetical world with no unmentioned factors, this establishes a sufficient condition: if \$Jyka\$ is true OR \$Yupt\$ is true, then \$Kwox\$ occurs (\$Kwox \textbackslash{}leftarrow Jyka \textbackslash{}lor Yupt\$).} \\
\\ 
\texttt{2.  **Analyze the Counterfactual Query**:} \\
\texttt{    *   **Observed World**: We observe an individual where \$Yupt = \textbackslash{}text\{False\}\$. The value of \$Jyka\$ in this observed world is irrelevant to the counterfactual question because we are explicitly asked to intervene on \$Jyka\$.} \\
\texttt{    *   **Intervention**: The question asks for the outcome if we set \$Jyka = \textbackslash{}text\{True\}\$ (replacing whatever its value was).} \\
\texttt{    *   **Resulting World**: In the counterfactual world, we have \$Jyka = \textbackslash{}text\{True\}\$ and \$Yupt = \textbackslash{}text\{False\}\$ (since we only intervened on \$Jyka\$).} \\
\\ 
\texttt{3.  **Determine the Outcome**:} \\
\texttt{    *   Since the condition for \$Kwox\$ is satisfied by the disjunction (\$Jyka \textbackslash{}lor Yupt\$), and we have established \$Jyka = \textbackslash{}text\{True\}\$ in the counterfactual scenario, the condition (\$True \textbackslash{}lor False\$) evaluates to True.} \\
\texttt{    *   Therefore, \$Kwox\$ must occur. The lack of information regarding the original state of \$Jyka\$ does not matter because the intervention forces \$Jyka\$ to be true, which is sufficient to trigger \$Kwox\$ regardless of \$Yupt\$'s status.} \\
            }
        }
    }
    & & \\ \\

    \theutterance \stepcounter{utterance}  
    & & & \multicolumn{2}{p{0.3\linewidth}}{
        \cellcolor[rgb]{0.9,0.9,0.9}{
            \makecell[{{p{\linewidth}}}]{
                \texttt{\tiny{[GM$|$GM]}}
                \texttt{No} \\
            }
        }
    }
    & & \\ \\

    \theutterance \stepcounter{utterance}  
    & & & \multicolumn{2}{p{0.3\linewidth}}{
        \cellcolor[rgb]{0.9,0.9,0.9}{
            \makecell[{{p{\linewidth}}}]{
                \texttt{\tiny{[GM$|$GM]}}
                \texttt{yes} \\
            }
        }
    }
    & & \\ \\

    \theutterance \stepcounter{utterance}  
    & & & \multicolumn{2}{p{0.3\linewidth}}{
        \cellcolor[rgb]{0.9,0.9,0.9}{
            \makecell[{{p{\linewidth}}}]{
                \texttt{\tiny{[GM$|$GM]}}
                \texttt{game\_result = LOSE} \\
            }
        }
    }
    & & \\ \\

\end{supertabular}
}

\end{document}
